%%% =======================================================================
%%% ISLS Extension Phase 7 v1.0.0
%%% C26: Crystallized Architecture Pattern Catalog
%%% Pre-Validated Skeletal Templates for Software Generation
%%% Author: Sebastian Klemm
%%% Date: 18 March 2026
%%% Status: Implementation Specification for Claude Code
%%% =======================================================================
\documentclass[11pt,a4paper,twoside]{article}

\usepackage[utf8]{inputenc}
\usepackage{amsmath,amssymb,amsthm}
\usepackage{booktabs,longtable,tabularx}
\usepackage{enumitem}
\usepackage{geometry}
\usepackage{hyperref}
\usepackage{cleveref}
\usepackage{xcolor}
\usepackage{fancyhdr}
\usepackage{listings}
\usepackage{titlesec}

\geometry{margin=2.5cm,top=3cm,bottom=3cm,headheight=14pt}

\theoremstyle{definition}
\newtheorem{definition}{Definition}[section]
\newtheorem{requirement}{Requirement}[section]
\newtheorem{invariant}{Invariant}[section]

\theoremstyle{plain}
\newtheorem{theorem}{Theorem}[section]

\theoremstyle{remark}
\newtheorem{remark}{Remark}[section]

\newcommand{\ISLS}{\textsf{ISLS}}

\definecolor{sec-color}{HTML}{1e3a5f}

\titleformat{\section}
  {\Large\bfseries\color{sec-color}}
  {\thesection}{1em}{}

\lstset{
  basicstyle=\ttfamily\small,
  breaklines=true,
  frame=single,
  backgroundcolor=\color{gray!5},
  rulecolor=\color{gray!30},
}

\pagestyle{fancy}
\fancyhf{}
\fancyhead[LE,RO]{\ISLS{} Extension Phase 7 v1.0.0}
\fancyhead[RE,LO]{Sebastian Klemm}
\fancyfoot[C]{\thepage}

\hypersetup{
  colorlinks=true,
  linkcolor=blue!60!black,
  pdftitle={ISLS Extension Phase 7 — Architecture Pattern Catalog},
  pdfauthor={Sebastian Klemm},
}

\begin{document}

\begin{titlepage}
\centering
\vspace*{2.5cm}
{\Huge\bfseries \ISLS{} Extension Phase~7\\[0.3cm]
v1.0.0}\\[1.5cm]
{\Large C26: Crystallized Architecture Pattern Catalog}\\[1cm]
{\large Pre-Validated Skeletal Templates\\
for Practical Software Generation}\\[0.5cm]
{\normalsize Normed structures so the Forge starts at 80\%,\\
not at 0\%.}\\[2cm]
{\large Sebastian Klemm}\\[0.5cm]
{\large 18 March 2026}\\[1.5cm]

\begin{abstract}
\noindent
This document specifies \texttt{isls-templates} (C26), a catalog of
pre-validated architecture patterns that are seeded into the Forge's
Pattern Memory at initialization. Each template is a complete
ArtifactIR tree (a CompositionTree from C24) encoding the standard
structure of a common software archetype: REST API, CLI tool, library,
microservice, database backend, WebSocket service, worker queue, and
more.

When the Forge receives a DecisionSpec, it first matches against the
template catalog by domain and archetype. If a template matches, the
Forge starts from that template's skeleton — not from scratch. The
Oracle (C25) then fills in the domain-specific logic, guided by the
template's component structure, interface contracts, and constraints.
The result is a complete, working project with correct architecture,
proper separation of concerns, and real code — not a flat string.

Templates are themselves crystals: content-addressed, validated
through the 8-gate cascade, and versioned in the registry. New
templates can be created by crystallizing successful forge outputs.
The catalog grows over time.

\medskip\noindent\textbf{Status:} Implementation Specification.\\
\textbf{Parent:} ISLS v1.0.0 + Extensions Phase 1--6.\\
\textbf{Depends on:} C22 (ArtifactIR), C23 (Forge), C24 (Compose),
  C25 (Oracle).
\end{abstract}
\end{titlepage}

\tableofcontents
\newpage


%%% =====================================================================
\part{Foundations}\label{part:foundations}

\section{The Problem}\label{sec:problem}

Without templates, the Forge receives a DecisionSpec like ``build a
REST API for user management'' and must decompose it from first
principles: what components does a REST API have? What interfaces
exist between routing, middleware, handlers, database, and error
handling? What are the standard patterns?

This means the PMHD drill, the Oracle, and the Compose engine must
rediscover well-known software architecture for every request. That
is wasteful: the architecture of a REST API is not a creative
problem. It is a solved problem with standard patterns. The creative
part is the \emph{domain-specific logic} inside those patterns.

Templates separate the solved structure from the unsolved content.

\section{What a Template Is}\label{sec:whatis}

\begin{definition}[Architecture Template]\label{def:template}
An architecture template $\tau$ is a pre-validated CompositionTree
(C24) where:
\begin{itemize}[nosep]
  \item Every atom has a complete ArtifactIR with components,
        interfaces, and constraints — but the component
        \texttt{content} fields contain \emph{structural skeletons}
        (type signatures, function signatures, trait definitions)
        rather than implementation logic.
  \item Every interface contract is fully specified (provider,
        consumer, type signatures, protocol).
  \item The composition passes CV1--CV6 (all interfaces resolved,
        no orphans, acyclic dependencies).
  \item The template itself is a Semantic Crystal in the archive,
        content-addressed and versioned.
\end{itemize}
\end{definition}

\begin{definition}[Template vs.\ Pattern]\label{def:tvsp}
A \emph{pattern} (in the Pattern Memory from C23/C25) is a validated
output from a previous forge run — it has full implementation code.
A \emph{template} is a structural skeleton without implementation —
it defines \emph{what} exists but not \emph{how} it works. Templates
are the starting point; patterns are the learned library.

\begin{tabularx}{\textwidth}{l X X}
\toprule
& \textbf{Template} & \textbf{Pattern} \\
\midrule
Content & Signatures, types, trait defs & Full implementation code \\
Source & Hand-crafted or distilled & Validated forge output \\
Purpose & Starting skeleton & Reuse of proven code \\
Matching & By archetype + domain & By 5D signature similarity \\
\bottomrule
\end{tabularx}
\end{definition}


%%% =====================================================================
\part{The Template Catalog}\label{part:catalog}

\section{Template Schema}\label{sec:schema}

\begin{lstlisting}[title={Template Structure}]
pub struct ArchitectureTemplate {
    pub id: Hash256,                    // content-addressed
    pub name: String,                   // "rest-api", "cli-tool", etc.
    pub version: String,                // semver
    pub domain: String,                 // "rust", "api", "workflow"
    pub archetype: Archetype,           // enum of known patterns
    pub description: String,
    pub tree: CompositionTree,          // from C24
    pub required_capabilities: Vec<String>, // what the Oracle must fill
    pub tags: Vec<String>,              // searchable tags
    pub crystal_id: Hash256,            // link to crystal in archive
}

pub enum Archetype {
    RestApi,
    CliTool,
    Library,
    Microservice,
    DatabaseBackend,
    WebSocketService,
    WorkerQueue,
    FullStackApp,
    DataPipeline,
    PluginSystem,
    Custom(String),
}

pub struct TemplateCatalog {
    templates: BTreeMap<String, ArchitectureTemplate>,
}

impl TemplateCatalog {
    pub fn load_defaults() -> Self;
    pub fn register(&mut self, template: ArchitectureTemplate);
    pub fn find_by_archetype(&self, arch: &Archetype) -> Vec<&ArchitectureTemplate>;
    pub fn find_by_tags(&self, tags: &[String]) -> Vec<&ArchitectureTemplate>;
    pub fn best_match(&self, spec: &DecisionSpec) -> Option<&ArchitectureTemplate>;
}
\end{lstlisting}


\section{Built-In Templates}\label{sec:builtin}

Phase 7 ships with 10 built-in templates. Each is described below
with its composition structure, components, and interfaces.

%%% -------------------------------------------------------------------
\subsection{T01: REST API Service}\label{sec:t01}

\begin{tabularx}{\textwidth}{l X}
\toprule
\textbf{Archetype} & \texttt{RestApi} \\
\textbf{Domain} & \texttt{rust} \\
\textbf{Description} & Axum-based REST API with layered architecture \\
\textbf{Tags} & rest, api, http, axum, crud, json \\
\bottomrule
\end{tabularx}

\medskip\noindent\textbf{Composition (3 molecules, 9 atoms):}

\begin{lstlisting}
System: rest-api-service
  |
  +-- Molecule: api-layer
  |     +-- Atom: router        (route definitions, method handlers)
  |     +-- Atom: middleware     (auth, logging, cors, rate-limit)
  |     +-- Atom: error_handler (error types, error responses)
  |
  +-- Molecule: domain-layer
  |     +-- Atom: models        (domain types, validation)
  |     +-- Atom: service       (business logic, trait definitions)
  |     +-- Atom: dto           (request/response types, serde)
  |
  +-- Molecule: infra-layer
        +-- Atom: database      (connection pool, migrations, queries)
        +-- Atom: config        (env vars, config struct, defaults)
        +-- Atom: main          (server startup, graceful shutdown)
\end{lstlisting}

\textbf{Key Interfaces:}
\begin{itemize}[nosep]
  \item \texttt{router} $\to$ \texttt{service}: handler calls service methods
  \item \texttt{service} $\to$ \texttt{database}: service uses DB trait
  \item \texttt{router} $\to$ \texttt{dto}: handlers deserialize/serialize DTOs
  \item \texttt{middleware} $\to$ \texttt{error\_handler}: middleware catches errors
  \item \texttt{main} $\to$ \texttt{config}: main reads config at startup
  \item \texttt{main} $\to$ \texttt{router}: main mounts router
\end{itemize}

\textbf{Dependencies (Cargo.toml skeleton):}
\begin{lstlisting}
axum, tokio, serde, serde_json, sqlx, tower-http,
tracing, tracing-subscriber, dotenvy, anyhow, thiserror
\end{lstlisting}

%%% -------------------------------------------------------------------
\subsection{T02: CLI Tool}\label{sec:t02}

\begin{tabularx}{\textwidth}{l X}
\toprule
\textbf{Archetype} & \texttt{CliTool} \\
\textbf{Domain} & \texttt{rust} \\
\textbf{Description} & Clap-based CLI with subcommands \\
\textbf{Tags} & cli, clap, command-line, tool \\
\bottomrule
\end{tabularx}

\medskip\noindent\textbf{Composition (2 molecules, 6 atoms):}

\begin{lstlisting}
System: cli-tool
  |
  +-- Molecule: interface
  |     +-- Atom: args          (Clap derive, subcommands, flags)
  |     +-- Atom: output        (table/json/plain formatting)
  |     +-- Atom: main          (entry point, error display)
  |
  +-- Molecule: core
        +-- Atom: config        (config file loading, defaults)
        +-- Atom: commands       (subcommand implementations)
        +-- Atom: errors        (error types, exit codes)
\end{lstlisting}

\textbf{Dependencies:} \texttt{clap, serde, serde\_json, toml, anyhow, thiserror}

%%% -------------------------------------------------------------------
\subsection{T03: Rust Library}\label{sec:t03}

\begin{tabularx}{\textwidth}{l X}
\toprule
\textbf{Archetype} & \texttt{Library} \\
\textbf{Domain} & \texttt{rust} \\
\textbf{Description} & Reusable Rust library with public API \\
\textbf{Tags} & library, crate, reusable, api \\
\bottomrule
\end{tabularx}

\medskip\noindent\textbf{Composition (2 molecules, 5 atoms):}

\begin{lstlisting}
System: rust-library
  |
  +-- Molecule: public-api
  |     +-- Atom: types         (public types, re-exports)
  |     +-- Atom: traits        (public traits, contracts)
  |     +-- Atom: lib           (lib.rs, module structure)
  |
  +-- Molecule: internals
        +-- Atom: impl          (trait implementations)
        +-- Atom: errors        (error types, From impls)
\end{lstlisting}

\textbf{Dependencies:} \texttt{serde (optional), thiserror}

%%% -------------------------------------------------------------------
\subsection{T04: Microservice}\label{sec:t04}

\begin{tabularx}{\textwidth}{l X}
\toprule
\textbf{Archetype} & \texttt{Microservice} \\
\textbf{Domain} & \texttt{rust} \\
\textbf{Description} & Production-ready microservice with health, metrics, graceful shutdown \\
\textbf{Tags} & microservice, health, metrics, docker, production \\
\bottomrule
\end{tabularx}

\medskip\noindent\textbf{Composition (4 molecules, 12 atoms):}

\begin{lstlisting}
System: microservice
  |
  +-- Molecule: api (from T01 api-layer, reused)
  +-- Molecule: domain (from T01 domain-layer, reused)
  +-- Molecule: infra (from T01 infra-layer, extended)
  |     +-- ... plus:
  |     +-- Atom: health        (liveness/readiness probes)
  |     +-- Atom: metrics       (prometheus metrics endpoint)
  |
  +-- Molecule: ops
        +-- Atom: dockerfile    (multi-stage Docker build)
        +-- Atom: ci            (GitHub Actions workflow)
        +-- Atom: env_template  (.env.example, secrets handling)
\end{lstlisting}

%%% -------------------------------------------------------------------
\subsection{T05: Database Backend}\label{sec:t05}

\begin{tabularx}{\textwidth}{l X}
\toprule
\textbf{Archetype} & \texttt{DatabaseBackend} \\
\textbf{Domain} & \texttt{rust} \\
\textbf{Description} & SQLx-based database layer with migrations \\
\textbf{Tags} & database, sqlx, sqlite, postgres, migrations, crud \\
\bottomrule
\end{tabularx}

\medskip\noindent\textbf{Composition (2 molecules, 6 atoms):}

\begin{lstlisting}
System: database-backend
  |
  +-- Molecule: schema
  |     +-- Atom: migrations    (SQL migration files, version tracking)
  |     +-- Atom: models        (Row types, FromRow derives)
  |     +-- Atom: queries       (CRUD operations, transactions)
  |
  +-- Molecule: connection
        +-- Atom: pool          (connection pool config)
        +-- Atom: testing       (test fixtures, in-memory DB)
        +-- Atom: errors        (database error types)
\end{lstlisting}

%%% -------------------------------------------------------------------
\subsection{T06: WebSocket Service}\label{sec:t06}

\begin{tabularx}{\textwidth}{l X}
\toprule
\textbf{Archetype} & \texttt{WebSocketService} \\
\textbf{Domain} & \texttt{rust} \\
\textbf{Description} & Axum WebSocket server with rooms and broadcast \\
\textbf{Tags} & websocket, realtime, broadcast, rooms, axum \\
\bottomrule
\end{tabularx}

\medskip\noindent\textbf{Composition (2 molecules, 6 atoms):}

\begin{lstlisting}
System: websocket-service
  |
  +-- Molecule: transport
  |     +-- Atom: ws_handler    (upgrade, message loop)
  |     +-- Atom: rooms         (room management, join/leave)
  |     +-- Atom: broadcast     (fan-out, backpressure)
  |
  +-- Molecule: protocol
        +-- Atom: messages      (message types, serde)
        +-- Atom: heartbeat     (ping/pong, timeout)
        +-- Atom: auth          (connection auth, token validation)
\end{lstlisting}

%%% -------------------------------------------------------------------
\subsection{T07: Worker / Job Queue}\label{sec:t07}

\begin{tabularx}{\textwidth}{l X}
\toprule
\textbf{Archetype} & \texttt{WorkerQueue} \\
\textbf{Domain} & \texttt{rust} \\
\textbf{Description} & Background job processor with retry and dead-letter \\
\textbf{Tags} & worker, queue, jobs, async, retry, background \\
\bottomrule
\end{tabularx}

\medskip\noindent\textbf{Composition (2 molecules, 6 atoms):}

\begin{lstlisting}
System: worker-queue
  |
  +-- Molecule: engine
  |     +-- Atom: dispatcher    (job dispatch, concurrency limits)
  |     +-- Atom: executor      (job execution, timeout)
  |     +-- Atom: retry         (retry policy, backoff, dead-letter)
  |
  +-- Molecule: jobs
        +-- Atom: job_trait     (Job trait, serialization)
        +-- Atom: registry     (job type registry)
        +-- Atom: store        (job persistence, status tracking)
\end{lstlisting}

%%% -------------------------------------------------------------------
\subsection{T08: Full-Stack Application}\label{sec:t08}

\begin{tabularx}{\textwidth}{l X}
\toprule
\textbf{Archetype} & \texttt{FullStackApp} \\
\textbf{Domain} & \texttt{rust} \\
\textbf{Description} & Axum backend + static frontend serving \\
\textbf{Tags} & fullstack, spa, static, axum, htmx \\
\bottomrule
\end{tabularx}

\medskip\noindent\textbf{Composition (4 molecules, 14 atoms):}

\begin{lstlisting}
System: fullstack-app
  |
  +-- Molecule: backend (reuses T04 microservice)
  +-- Molecule: frontend
  |     +-- Atom: pages         (HTML templates / components)
  |     +-- Atom: assets        (CSS, JS, static files)
  |     +-- Atom: build         (build script, bundling)
  |
  +-- Molecule: shared
  |     +-- Atom: types         (shared types between FE/BE)
  |     +-- Atom: validation    (shared validation rules)
  |
  +-- Molecule: deployment
        +-- Atom: dockerfile
        +-- Atom: nginx_conf    (reverse proxy config)
\end{lstlisting}

%%% -------------------------------------------------------------------
\subsection{T09: Data Pipeline}\label{sec:t09}

\begin{tabularx}{\textwidth}{l X}
\toprule
\textbf{Archetype} & \texttt{DataPipeline} \\
\textbf{Domain} & \texttt{rust} \\
\textbf{Description} & ETL/streaming pipeline with stages \\
\textbf{Tags} & pipeline, etl, streaming, transform, data \\
\bottomrule
\end{tabularx}

\medskip\noindent\textbf{Composition (3 molecules, 8 atoms):}

\begin{lstlisting}
System: data-pipeline
  |
  +-- Molecule: ingestion
  |     +-- Atom: source        (source trait, file/http/stdin)
  |     +-- Atom: parser        (format detection, deserialization)
  |
  +-- Molecule: processing
  |     +-- Atom: transform     (map, filter, aggregate, window)
  |     +-- Atom: validate      (schema validation, quality checks)
  |     +-- Atom: enrich        (join, lookup, derive)
  |
  +-- Molecule: output
        +-- Atom: sink          (sink trait, file/db/http)
        +-- Atom: checkpoint    (progress tracking, resume)
        +-- Atom: metrics       (throughput, latency, error rates)
\end{lstlisting}

%%% -------------------------------------------------------------------
\subsection{T10: Plugin System}\label{sec:t10}

\begin{tabularx}{\textwidth}{l X}
\toprule
\textbf{Archetype} & \texttt{PluginSystem} \\
\textbf{Domain} & \texttt{rust} \\
\textbf{Description} & Extensible host with dynamic plugin loading \\
\textbf{Tags} & plugin, extensible, dynamic, trait-object, registry \\
\bottomrule
\end{tabularx}

\medskip\noindent\textbf{Composition (2 molecules, 6 atoms):}

\begin{lstlisting}
System: plugin-system
  |
  +-- Molecule: host
  |     +-- Atom: plugin_trait  (Plugin trait, lifecycle hooks)
  |     +-- Atom: registry      (plugin discovery, registration)
  |     +-- Atom: loader        (dynamic loading or static dispatch)
  |
  +-- Molecule: runtime
        +-- Atom: sandbox       (resource limits per plugin)
        +-- Atom: events        (event bus, plugin communication)
        +-- Atom: config        (per-plugin configuration)
\end{lstlisting}


%%% =====================================================================
\part{Template Matching and Adaptation}\label{part:matching}

\section{How Templates Are Selected}\label{sec:select}

\begin{definition}[Template Selection Algorithm]\label{def:select}
When the Forge receives a DecisionSpec, the template catalog is
queried:
\begin{enumerate}[nosep]
  \item \textbf{Archetype match}: if \texttt{spec.archetype} is
        set, filter templates by archetype.
  \item \textbf{Tag match}: compute tag overlap between spec tags
        and template tags. Rank by overlap ratio.
  \item \textbf{Domain match}: filter by \texttt{spec.domain}.
  \item \textbf{Keyword match}: if no archetype is specified, match
        spec intent keywords against template descriptions and tags.
  \item \textbf{Best match}: return the template with highest
        combined score, or \texttt{None} if no template scores
        above threshold.
\end{enumerate}
\end{definition}

\begin{lstlisting}[title={Template Matching in ForgeEngine}]
impl ForgeEngine {
    pub fn forge(&mut self, spec: DecisionSpec) -> Result<ForgeResult> {
        // Step 1: Check template catalog
        let template = self.catalog.best_match(&spec);

        if let Some(tmpl) = template {
            // Step 2: Clone template's CompositionTree
            let mut tree = tmpl.tree.clone();

            // Step 3: Adapt tree to spec
            //   - Rename components to match spec intent
            //   - Inject spec constraints into atom sub-specs
            //   - Mark atoms that need Oracle fill
            self.adapt_template(&mut tree, &spec);

            // Step 4: For each atom that needs fill:
            //   Oracle synthesizes the implementation
            self.fill_atoms(&mut tree, &spec);

            // Step 5: Compose upward (CV1-CV6 + 8-gate)
            self.compose_upward(&mut tree)
        } else {
            // No template match: full decomposition path (C24)
            self.forge_from_scratch(spec)
        }
    }
}
\end{lstlisting}


\section{Atom Fill Strategies}\label{sec:fill}

\begin{definition}[Fill Strategy]\label{def:fill}
Each atom in a template tree has a fill strategy:
\begin{itemize}[nosep]
  \item \textbf{Oracle}: the Oracle (C25) generates the
        implementation. Used for domain-specific logic (handlers,
        business rules, queries).
  \item \textbf{Pattern}: the Pattern Memory provides a proven
        implementation from a prior run. Used for standard
        components (error types, config loading, connection pools).
  \item \textbf{Static}: the template itself contains a complete,
        hardcoded implementation. Used for boilerplate (Cargo.toml,
        Dockerfile, CI config).
  \item \textbf{Derive}: the implementation is derived from other
        atoms' types (e.g., DTO derives from model types).
\end{itemize}
\begin{lstlisting}
pub enum FillStrategy {
    Oracle,
    Pattern,
    Static { content: String },
    Derive { source_atom: String, transform: String },
}
\end{lstlisting}
\end{definition}

\begin{definition}[Atom Metadata in Template]\label{def:atommeta}
Each atom in a template carries fill guidance:
\begin{lstlisting}
pub struct TemplateAtom {
    pub name: String,
    pub fill_strategy: FillStrategy,
    pub skeleton: String,           // type signatures, trait defs
    pub oracle_hint: Option<String>, // hint for Oracle prompt
    pub constraints: Vec<String>,    // hard constraints
    pub test_requirements: Vec<String>, // what tests should exist
}
\end{lstlisting}
\end{definition}


%%% =====================================================================
\part{Template Lifecycle}\label{part:lifecycle}

\section{Creating New Templates}\label{sec:create}

Templates can be created in three ways:

\begin{definition}[Manual Template Creation]\label{def:manual}
A developer defines a template by hand:
\begin{lstlisting}
isls template create --name "my-pattern" --archetype rest-api \
  --structure structure.json --domain rust
\end{lstlisting}
The structure file defines the composition tree, atoms, interfaces,
and fill strategies.
\end{definition}

\begin{definition}[Template Distillation]\label{def:distill}
A successful forge output can be distilled into a template:
\begin{lstlisting}
isls template distill --crystal <crystal_id> --name "proven-api"
\end{lstlisting}
This strips the implementation code from the crystal's ArtifactIR
tree, retains the structure (components, interfaces, constraints),
and saves the result as a new template. The implementation code
remains in the Pattern Memory for reuse.
\end{definition}

\begin{definition}[Template Composition]\label{def:compose}
Templates can be composed from sub-templates:
\begin{lstlisting}
isls template compose --name "microservice" \
  --include rest-api --include database-backend \
  --include health-metrics
\end{lstlisting}
This merges the composition trees of multiple templates, resolves
shared interfaces, and produces a new composite template.
\end{definition}


\section{Template Versioning}\label{sec:version}

\begin{requirement}[Template Versioning]\label{req:version}
Templates use semantic versioning. When a template is updated, the
old version remains in the catalog (templates are crystals --- they
are immutable). The catalog tracks which version is ``active'' for
each archetype.
\end{requirement}


%%% =====================================================================
\part{Configuration and CLI}\label{part:config}

\section{Configuration}\label{sec:config}

\begin{lstlisting}[title={Template Configuration}]
templates:
  catalog_dir: "~/.isls/templates/"    # template crystal storage
  auto_match: true                     # try templates before scratch
  match_threshold: 0.3                 # minimum tag overlap
  distill_on_success: true             # auto-distill successful forges
  active_templates:
    rest-api: "v1.0.0"
    cli-tool: "v1.0.0"
    # ... etc.
\end{lstlisting}

\section{CLI}\label{sec:cli}

\begin{lstlisting}
# List available templates
isls template list
# -> T01 rest-api v1.0.0 (9 atoms, 3 molecules)
# -> T02 cli-tool v1.0.0 (6 atoms, 2 molecules)
# -> ...

# Show template structure
isls template show rest-api
# -> Composition tree with atoms, interfaces, fill strategies

# Create from structure file
isls template create --name custom --structure custom.json

# Distill from a forge result
isls template distill --crystal <id> --name proven-pattern

# Compose templates
isls template compose --name mega --include rest-api --include websocket

# Forge using a specific template
isls forge --spec intent.json --template rest-api

# Forge with auto-template (default)
isls forge --spec intent.json --domain rust
# Template auto-selected based on intent keywords
\end{lstlisting}

\section{Gateway}\label{sec:gateway}

\begin{lstlisting}[title={New Gateway Endpoints}]
GET  /templates              List available templates
GET  /templates/{name}       Show template structure
POST /templates              Create new template
POST /templates/distill      Distill template from crystal
POST /forge                  Now accepts optional template_name field
\end{lstlisting}


%%% =====================================================================
\part{Acceptance Tests}\label{part:tests}

\begin{longtable}{l l p{5.5cm}}
\toprule
\textbf{ID} & \textbf{Name} & \textbf{Procedure} \\
\midrule
\endfirsthead \midrule \endhead \bottomrule \endfoot
AT-T1 & Catalog load & Load default catalog; verify 10 templates
  present. \\
AT-T2 & Archetype match & Query for RestApi archetype; verify T01
  returned. \\
AT-T3 & Tag match & Query with tags [``websocket'', ``realtime''];
  verify T06 ranked highest. \\
AT-T4 & Template structure & Load T01; verify 9 atoms, 3 molecules,
  6 interfaces. \\
AT-T5 & Template is crystal & Verify each built-in template has a
  valid crystal ID and passes 8-gate. \\
AT-T6 & Forge with template & Forge ``user management API'' with T01;
  verify output has router, service, database atoms. \\
AT-T7 & Auto-match & Forge with intent containing ``REST API'';
  verify T01 auto-selected. \\
AT-T8 & Distillation & Forge a project; distill result; verify new
  template in catalog with same structure but no impl code. \\
AT-T9 & Template composition & Compose T01 + T05; verify merged tree
  with shared database atom. \\
AT-T10 & Fill strategies & Verify T01 has: Oracle for handlers,
  Static for Cargo.toml, Derive for DTOs. \\
AT-T11 & No template fallback & Forge with exotic intent that matches
  no template; verify fallback to full decomposition. \\
AT-T12 & Template versioning & Create v2 of a template; verify both
  v1 and v2 exist; verify active version used by default. \\
\end{longtable}


%%% =====================================================================
\part{Handoff}\label{part:handoff}

\section{Test Count}

\begin{tabular}{l r r}
\toprule
\textbf{Phase} & \textbf{New Tests} & \textbf{Cumulative} \\
\midrule
Phase 1--6 + Genesis & 253 & 253 \\
Phase 7 (C26) & 12 & $\ge$ 265 \\
\bottomrule
\end{tabular}

\section{Completion Criteria}

\begin{enumerate}[nosep]
  \item 10 built-in templates loaded at startup.
  \item Template matching selects correct template by archetype, tags,
        and keywords.
  \item Forge with template produces structured, multi-file output.
  \item Each template is a valid crystal.
  \item Distillation creates new templates from forge results.
  \item Template composition merges trees with shared interfaces.
  \item Auto-match works for common intents.
  \item Fallback to full decomposition when no template matches.
  \item 12 acceptance tests pass.
  \item Total tests $\ge 265$, zero warnings.
  \item \textbf{26 Crates. The Forge starts at 80\%, not at 0\%.}
\end{enumerate}


\vfill
\begin{center}
\rule{0.5\textwidth}{0.4pt}\\[0.5cm]
{\small Generated for \ISLS{} Extension Phase 7 v1.0.0.\\
10 archetypes. Normed structures. Domain-specific fill.\\
The solved parts are solved. The creative parts are where the Oracle works.}
\end{center}

\end{document}
